\section{Appendix}\label{sec:appendix}

\subsection{Is the signal in the features? The ranking probe}\label{app:probe}
The group task asks which member does not belong. We test whether that is
answerable from the aligned features alone, without a language model in the
loop. For each group we rank its members by one quantity and ask where a
qualifying star lands. Two readouts are compared on exactly the same groups:
the true catalog value, and a cross-fit ridge regression that predicts the same
quantity from the aligned feature (5-fold, so no star is scored by a model that
saw it). If the probe tracks the catalog, the signal survives the alignment.

Two details matter. First, only the three continuous axes can be ranked by a
scalar; binarity and cluster membership are categorical and are excluded
($103$ of $259$ anomalies). Second, the composer designates one intruder but
leaves the other members unconstrained on that axis, so a group often contains
several stars carrying the same label --- $60.5\%$ of anomalies on the full
fold, and $81.7\%$ on dwarf/giant. Naming any of them is a true statement about
the group, so a readout is credited for the best-placed qualifying star rather
than only the designated one. ``Extreme'' is distance from the group median in
units of its median absolute deviation, which does not reveal which direction
the answer lies in.

\begin{table}[h]\centering\small
\caption{\textbf{Ranking the odd member from the features alone.} $n{=}156$
continuous-axis anomalies over $129$ groups. \emph{rank-1} = a qualifying star
is the single most extreme member; \emph{top-3} = one is among the three most
extreme. The probe reads only what the language model reads.}
\label{tab:probe}
\begin{tabular}{lcc}
\toprule
readout & rank-1 & top-3 \\
\midrule
catalog values (oracle)      & 0.660 & \textbf{0.942} \\
probe, shared subspace (64-d) & 0.660 & 0.840 \\
probe, full feature (1024-d)  & \textbf{0.686} & 0.891 \\
\midrule
\multicolumn{3}{l}{\emph{by axis, full feature vs.\ catalog}}\\
dwarf / giant     & 0.911 \emph{vs} 0.857 & 1.000 \emph{vs} 0.929 \\
hot / cool        & 0.755 \emph{vs} 0.694 & 0.939 \emph{vs} 0.959 \\
metal rich / poor & 0.373 \emph{vs} 0.412 & 0.725 \emph{vs} 0.941 \\
\bottomrule
\end{tabular}
\end{table}

The probe matches the catalog at rank-1 and keeps $89$--$95\%$ of it at top-3;
on the gravity axis it is ahead of the catalog. The information the group task
needs is in the features.

The catalog row is not a ceiling on the task. It bounds what a \emph{single
ranked quantity} can do, and it falls short of $1.0$ because the label is a
threshold test (``$[\mathrm{Fe/H}] > 0.3$'') while a ranking is a distance test:
a star can sit far from the group median without crossing the threshold. A
readout that models the group jointly --- the head of
\Cref{tab:headboth} --- is not bound by it.

Scoring the narrators by the same rule, crediting any equally-qualifying star,
raises both arms and leaves the head's lead intact: $+0.361 \to +0.318$
($n{=}1{,}203$ groups, $p{=}6\times10^{-85}$).

\subsection{The regime control: same prose, same LoRA, no token}\label{app:2x2}
The head arms are trained with LoRA on a frozen base while the no-head arms are
full fine-tunes, so ``head vs.\ no head'' is also ``LoRA vs.\ full fine-tune''.
We closed that on the galaxy track, where one cell costs $\sim$10 GPU-hours: a
narrator trained on the \emph{same} concordance prose, with the \emph{same} LoRA
rank and schedule as the head arm, and no cross-modal token. It is one variable
away from the head arm.

\begin{table}[h]\centering\small
\caption{\textbf{The complete $2\times2$}, galaxy cross-modal task, minimum-AUC
over the two star-formation strata ($0.50$ = chance; brackets are $95\%$
intervals on the minority stratum). A score above $0.5$ in \emph{both} strata
detects the interaction; below it, the arm is reading morphology alone.}
\label{tab:2x2}
\begin{tabular}{lccc}
\toprule
 & full fine-tune & LoRA & read \\
\midrule
no token  & 0.442 [0.409,0.477] & 0.498 [0.462,0.535] & morphology-only \\
$+$ token & ---                 & \textbf{0.628} [0.593,0.665] & interaction \\
\bottomrule
\end{tabular}
\end{table}

Neither training regime produces an interaction detector on its own: both
no-token arms sit at chance on the quiescent stratum, and the LoRA arm's
interval contains $0.5$. The full fine-tune is in fact the more confident
shortcut learner --- it posts the highest \emph{marginal} AUC in the study
($0.885$) while scoring $0.442$ on the minority stratum, which is why we never
quote marginal AUC here. Only the arm with the token is above chance in both
strata. We take this as evidence, by analogy, for the stellar track as well,
where the equivalent control would cost $\sim$700--1{,}000 GPU-hours and was
not run.
